\subsection{The Gaussian seed, its two couplings, and the complementary evolution}
\label{sec:short-gaussian-seed}

Use $u=\nabla^\perp\psi$, $\operatorname{curl}u=\Delta\psi$ and put
\[
 x=r^2/4,\quad \Omega=e^{-x},\quad
 V_0(y)=\Omega y^\perp,\quad E_0=(2-2x)e^{-x},\quad
 k=E_0'/r=(x-2)e^{-x},\quad G=(4\pi)^{-1}e^{-x}.
\]
The Euler linearization is
\[
 \mathscr Lf=-\Omega\partial_\vartheta f
       -(\nabla^\perp\Delta^{-1}f)\cdot\nabla E_0.
\]
For $j\ge0$, $M>1$, let $H^j_M$ be the space with norm
$\sum_{|\alpha|\le j}\|\langle y\rangle^M\partial^\alpha f\|_2$ and let
$\mathcal X_{j,M}=\{f\in H^j_M:\int f=0\}$.
In angular sector two use
\[
 X_G=L^2((0,\infty),G^{-1}r\,dr)e^{2i\vartheta},\qquad
 L_0=-2i\Omega+2ikR,\quad R=\Delta_2^{-1},
\]
where the pairing is linear in its first argument and omits $2\pi$.
For a general planar function, $f_2(r)=(2\pi)^{-1}\int_0^{2\pi}
f(r,\vartheta)e^{-2i\vartheta}d\vartheta$.
The regular decaying Green inverse in angular sector $n\ge1$ is
\begin{equation}\label{eq:short-seed-green}
 R_nw=-\frac1{2n}\left[
 r^{-n}\int_0^r s^{n+1}w(s)\,ds+
 r^n\int_r^\infty s^{1-n}w(s)\,ds\right].
\end{equation}

\begin{proposition}\label{prop:short-gaussian-seed}
On every $\mathcal X_{j,M}$ the entire spectrum of $\mathscr L$ in
$\{\Re z>0\}$ is one algebraically simple conjugate pair
$\lambda=\gamma+i\omega$, $\bar\lambda$, where $\gamma>0$ and
$-2<\omega<0$. Its eigenvectors are $q=h(r)e^{2i\vartheta}$ and
$\bar q$. Every derivative of $q$ is bounded by a polynomial times
$e^{-r^2/4}$, $\int q=0$, its velocity belongs to $L^2$, and
\[
 M_q:=\int_0^\infty r^3h(r)\,dr\ne0.
\]
The same $\lambda$ is the sole growing spectral value of $L_0$ on $X_G$.
Its projection is $P=q\ell$, with
\begin{equation}\label{eq:short-seed-functional}
 \ell(f)=\frac1{J_q}\int_0^\infty
       \frac{h(r)f_2(r)}{k(r)}r\,dr,\qquad
 J_q=\int_0^\infty\frac{h(r)^2}{k(r)}r\,dr\ne0.
\end{equation}
The quotients at $k=0$ are removable. In particular, the harmonic-strain
input $H=r^2k(r)e^{2i\vartheta}$ satisfies $\ell(H)=M_q/J_q\ne0$.

With $\mathscr P$ the projection onto the pair and $Q=I-P$,
\begin{equation}\label{eq:short-seed-complement}
 \|e^{t\mathscr L}(I-\mathscr P)\|_{H^j_M\to H^j_M}
       \le C_{j,M,\eta}e^{\eta t},\qquad
 \|e^{t(L_0-\lambda)}Q\|_{X_G\to X_G}
       \le C_\eta e^{-(\gamma-\eta)t}
\end{equation}
for every $\eta>0$ and $t\ge0$. Finally
$D=\Delta+\frac12y\cdot\nabla+2$ is selfadjoint and nonpositive on
$X_G$, with spectrum $\{0,-1,-2,\ldots\}$. Both $q$ and the normalized
Gaussian adjoint vector $q^\dagger$ belong to every $\mathcal D(D^m)$;
thus $DP$ and $PD$ extend boundedly to $X_G$.
\end{proposition}

\begin{proof}
\emph{One negative direction.}
Set $a_*=e^{-2}$, $t=\log r$, and
\[
 v=-\frac{k}{\Omega-a_*}
       =\frac{x-2}{e^{x-2}-1}>0,\qquad
 W=r^2v=\frac{4x(x-2)}{e^{x-2}-1},\qquad
 \mathcal H_n=-\partial_t^2+n^2-W.
\]
At $x=2$, define $v=1$ and $W=8$ by continuity. The bounded function $W$ tends to
zero at both ends, so $\mathcal H_n$ is selfadjoint on $H^2(\mathbb R)$
and has essential spectrum $[n^2,\infty)$: multiplication by $W$ is
relatively compact by truncation and Rellich's theorem. Also $\sup W<9$: for $y=x-2$, the function
$F(y)=9(e^y-1)-8y-4y^2$ satisfies $F(0)=0$ and
$F'(y)\ge-8\log(8/9)>0$, so $9-W=F(y)/(e^y-1)>0$;
at $y=0$ the inequality follows from $W=8$.

There is a negative trial in sector two. For
$\phi=r^2e^{-r^2/8}$ direct integration gives
\[
 \int_0^\infty(|\phi'|^2+4r^{-2}|\phi|^2-v|\phi|^2)r\,dr
 =48-32\int_0^\infty\frac{x-2}{e^{x-2}-1}x^2e^{-x}\,dx.
\]
For all real $y$,
\begin{equation}\label{eq:short-seed-rational}
 \frac y{e^y-1}\ge1-\frac y2+\frac{5y^2}{60+y^2}.
\end{equation}
Indeed, with $z=y/2$, this is equivalent to
\[
 \frac{(z^3+15z)\cosh z-(6z^2+15)\sinh z}{\sinh z}\ge0.
\]
At $z=0$ take the continuous extension. The numerator has coefficients
$(2m)(2m-2)(2m-4)/(2m+1)!\ge0$ at $z^{2m+1}$.
For $d\rho=x^2e^{-x}dx$, $y=x-2$, and $p=y(16-y)/930$,
completion of the square and $\int_0^\infty x^ke^{-x}dx=k!$ give
\[
 \int\frac{y^2}{60+y^2}\,d\rho
 \ge\int\bigl(2yp-(60+y^2)p^2\bigr)\,d\rho=\frac{16}{155}.
\]
Thus the trial form is at most $-16/31$.

There cannot be a second nonpositive direction. The function
$h_*(t)=e^t(\Omega-a_*)$ has one simple zero
$t_*=\log\sqrt8$ and $h_*''/h_*=1-W$. Ground-state factorization
on the two half-lines gives
\[
 \int(|F'|^2+(1-W)|F|^2)\,dt
    =\int|F'-(h_*'/h_*)F|^2\,dt\ge0
 \quad\text{when }F(t_*)=0.
\]
Prove it first for compact support away from $t_*$ and then use
$H^1_0$ density. Hence
$\langle F,\mathcal H_nF\rangle\ge(n^2-1)\|F\|_2^2$ on a
codimension-one subspace. Min--max implies that $\mathcal H_2$ has
exactly one negative eigenvalue, counting multiplicity, and no kernel.
For $n\ge3$ the bound $W<9$ directly gives
\begin{equation}\label{eq:short-seed-higher}
 \langle F,\mathcal H_nF\rangle\ge
 (1-\|W\|_\infty/n^2)(\|F'\|_2^2+n^2\|F\|_2^2).
\end{equation}

\emph{A bounded Hamiltonian realization and its neutral modes.}
On $Y=L^2((0,\infty),r\,dr)$ set
\[
 K_n=\sqrt v(-\Delta_n)^{-1}\sqrt v,\quad B_n=I-K_n,
 \quad J=\Omega-a_*,\quad \widetilde T_n=a_*+JB_n.
\]
The kernel of $K_n$, relative to $s\,ds$, is
$\frac1{2n}\sqrt{v(r)v(s)}(\min(r,s)/\max(r,s))^n$.
Since $\int rv(r)dr<\infty$, it is Hilbert--Schmidt; Green's identity
makes it nonnegative and selfadjoint. More explicitly, on the energy
space $\mathcal E_n$ with norm
$\int(|\phi'|^2+n^2r^{-2}|\phi|^2)r\,dr$, the compact map
$S\phi=\sqrt v\phi$ satisfies $SS^*=K_n$.
The nonzero eigenvalues of $SS^*$ and $S^*S$ agree with multiplicity.
The preceding form computation therefore says that $B_2$ is invertible
with one complex negative direction, whereas $B_n$ is strictly positive
for $n\ge3$. The change $w=\sqrt v f$ intertwines
$T_n=\Omega-k\Delta_n^{-1}$ with $\widetilde T_n$.

If $\widetilde T_nf=cf$ and $\psi=S^*f$, then
$F(t)=\psi(e^t)\in H^2(\mathbb R)$ and
\begin{equation}\label{eq:short-seed-rayleigh}
 (\Omega-c)(F''-n^2F)-AF=0,\qquad
 A=r^2k=\partial_t^2\Omega+2\partial_t\Omega.
\end{equation}
Indeed the Green source $\rho=r^2\sqrt v f$ has squared logarithmic
$L^2$ norm $\int W|f|^2r\,dr$, and
$-F''+n^2F=\rho$. Also $F=0$ forces $f=0$.

The operator $\widetilde T_2$ has no real eigenvalues. At $c=a_*$ the equation
$JB_2f=0$ gives $f=0$, since $J$ vanishes at one point and $B_2$ is
invertible. At $c=0$, \eqref{eq:short-seed-rayleigh} becomes
$-F''+4(x-1)^2F=0$, whose cutoff energy identity gives $F=0$.
For the other real $c$, put $h_c=e^t(\Omega-c)$; then
$h_c''/h_c=1+A/(\Omega-c)$ and the equation is
$-F''+(3+h_c''/h_c)F=0$.
When $0<c<1$, $c\ne a_*$, its only critical point $t_c$ has
$A(t_c)\ne0$. The equation and $F''-4F\in L^2$ force
$F(t_c)=0$, since otherwise $AF/(\Omega-c)$ has a nonintegrable
simple pole. As $F\in H^2$, $F=O(|t-t_c|)$.
For $c<0$ or $c\ge1$ there is no finite zero. In every case,
factorization on the intervals avoiding $t_c$ gives
\[
 0=\int\left(|F'-(h_c'/h_c)F|^2+3|F|^2\right)dt.
\]
The critical boundary term is $O(|t-t_c|)$; the terms at infinity
vanish because $F,F'\to0$ and $h_c'/h_c$ is bounded there
(its limit at $-\infty$ is three if $c=1$). Thus $F=0$.

Regard $Y$ as a real Hilbert space. The operator
$\mathcal A=-2iJB_2$ has bounded skew-adjoint factor $-2iJ$
and bounded selfadjoint energy $B_2$. Its negative index is two
over the reals, its kernel is zero, and its positive space is
uniformly positive since $B_2$ is an invertible compact perturbation
of identity. These verify the hypotheses of the index theorem of Lin and
Zeng~\cite{lin-zeng2021}, which gives
\[
 k_r+2k_c+2k_i^{\le0}+k_0^{\le0}=2.
\]
The last two counts vanish by the preceding absence of imaginary
eigenvectors; continuous spectrum creates no additional term in this
theorem. The complexification of the real operator is unitarily
equivalent to $\mathcal A\oplus\overline{\mathcal A}$ on $Y\oplus Y$,
so its total unstable algebraic multiplicity is twice that of
$\mathcal A$ on complex $Y$. Hence it is one on $Y$.
Equivalently, $\widetilde T_2$ has exactly one upper-half-plane
eigenvalue $c$, algebraically simple.

\emph{Frequency, Gaussian regularity, and the quadrupole.}
For this nonreal $c$, substitute $F=(\Omega-c)\Phi$ into
\eqref{eq:short-seed-rayleigh} and integrate against $\bar\Phi$:
\[
 \int(\Omega-c)^2\Theta\,dt=0,\qquad
 \Theta=|\Phi'|^2+4|\Phi|^2-2\Re(\Phi'\bar\Phi)
          \ge3|\Phi|^2.
\]
Here $\Phi\in H^2$, so integration by parts is valid.
The imaginary and real parts imply respectively
$\int\Omega\Theta=(\Re c)\int\Theta$ and
$\int\Omega^2\Theta=|c|^2\int\Theta$.
Since $0<\Omega<1$, $|c-1/2|<1/2$. Thus
$\lambda=-2ic$ has the asserted growth and frequency.

Write $\phi(r)=F(\log r)$. The coefficient
$B=A/(\Omega-c)$ in $F''-4F=BF$ is $O(e^{2t})$ at $-\infty$
and decreases faster than every exponential at $+\infty$, with
the same differentiated bounds. At $+\infty$ the Volterra equation
\[
 F_+(t)=e^{-2t}+\frac12\int_t^\infty
                   \sinh(2(s-t))B(s)F_+(s)\,ds
\]
is a contraction on $\sup_{t\ge T}e^{2t}|F(t)|$ for large $T$.
It gives $F_+=e^{-2t}(1+o(1))$ with differentiated asymptotics;
reduction of order supplies an independent solution growing as
$e^{2t}$. The analogous forward Volterra construction at $-\infty$
selects $e^{2t}(1+O(e^{2t}))$. Thus every nonzero $H^1$ solution has
\[
 \phi(r)=B_0r^2(1+O(r^2))\ (r\downarrow0),\qquad
 \phi(r)=C_\infty r^{-2}(1+o(1))\ (r\to\infty),
 \qquad B_0C_\infty\ne0.
\]
The Cartesian streamfunction $\Psi=\phi e^{2i\vartheta}$ is in $L^2$.
These endpoint estimates remove the inner-circle boundary terms in
Green's identity, so
$\Delta\Psi=[k/(\Omega-c)]\Psi$ on the whole plane.
The coefficient is smooth with bounded derivatives. The estimate
$\|\Psi\|_{H^2}\le C(\|\Delta\Psi\|_2+\|\Psi\|_2)$,
differentiated and iterated, gives all Sobolev norms and bounded
Cartesian derivatives. Hence
$q=\Delta\Psi=h e^{2i\vartheta}$, with
$h=k\phi/(\Omega-c)$, has the claimed Gaussian bounds.
Its velocity is $O(r)$ at zero and $O(r^{-3})$ at infinity.
Finally
\begin{equation}\label{eq:short-seed-quadrupole}
 M_q=\int_0^\infty r^3\Delta_2\phi\,dr
       =[r^3\phi'-2r^2\phi]_0^\infty=-4C_\infty\ne0.
\end{equation}
Its circulation vanishes by angular integration.

\emph{Excluding other angular sectors and transferring rootspaces.}
For $n\ge3$, $B_n>0$ and
$\widetilde T_n^*B_n=B_n\widetilde T_n$ imply every phase-speed
eigenvalue is real. For $n=1$, $\phi_1=r(\Omega-c)$ is an exact
solution for every nonreal $c$. Its independent solution is
$\phi_1(r)\int^r ds/(s\phi_1(s)^2)
=-[2(1-c)r]^{-1}+O(r|\log r|)$ at zero, and is not locally
square integrable. The regular solution $\phi_1$ grows as $-cr$
at infinity, excluding a regular decaying mode. Radial vorticities
are stationary. Negative sectors are the conjugates of positive ones.

This classification applies to the actual weighted spaces without
claiming a bounded similarity $w\mapsto w/\sqrt v$ there.
Cauchy--Schwarz in \eqref{eq:short-seed-green} gives, for $r\ge2$,
\[
 |R_nw(r)|\le C_{n,M}\|w\|_{L^2_M}
       \left(r^{1-M}+r^{-n}(1+\log r)^{1/2}\right);
\]
at zero it is $O(r)$, with an extra $|\log r|^{1/2}$ only for $n=1$.
For a nonreal rootvector the equation
$w=kR_nw/(\Omega-c)$, and for each subsequent Jordan vector the
same equation with its preceding vector divided by $\Omega-c$,
show inductively that $w/\sqrt v\in Y$.
Indeed $k/\sqrt v$ is Gaussian at infinity and bounded at zero.
The Green inverse agrees with $-S^*(w/\sqrt v)$: their difference
is $ar^n+br^{-n}$, and both vanish at the two endpoints.
Thus whole Jordan chains transfer injectively to $Y$.
Local elliptic regularity selects precisely the regular solutions
used above. In $X_G$ membership in $Y$ is immediate from $v\ge cG$.
Conversely the constructed Gaussian eigenmode belongs to all these
spaces, so the growing eigenvalue and its algebraic multiplicity agree.

\emph{Generators, projections, and complementary bounds.}
The free transport is
$T(t)f(r,\vartheta)=f(r,\vartheta-\Omega(r)t)$.
The chain rule gives $\|T(t)\|_{H^j_M\to H^j_M}\le C_j(1+|t|)^j$.
Density and differentiation of the flow give its generator domain
$\{f:\Omega\partial_\vartheta f\in H^j_M\}$.
For the converse domain inclusion use the distributional identity
$T(t)f-f=-\int_0^tT(s)\Omega\partial_\vartheta f\,ds$ and strong continuity.
The interaction $(\nabla^\perp\Delta^{-1}f)\cdot\nabla E_0$ is
compact on $H^j_M$: splitting the Biot--Savart kernel at distance
one controls its undifferentiated near part in $L^2$ and far part
in $L^\infty$ by $\|f\|_{L^2}+\|f\|_1$; derivatives have Fourier
order at most $j$. Multiplication by the Schwartz coefficient
$\nabla E_0$ therefore maps $H^j_M$ into $H^{j+1}_M$.
Truncation to a disk, Rellich compactness, and the vanishing weighted
coefficient tails prove compactness on $H^j_M$.
The bounded perturbation theorem now realizes $\mathscr L$ on the
stated domain. Both terms have zero integral, so $\mathcal X_{j,M}$
is invariant.

The Duhamel difference $e^{t\mathscr L}-T(t)$ is compact for each
$t>0$: its integrand consists of the compact interaction between
strongly continuous evolutions, and is norm continuous because the
adjoint evolutions are strongly continuous on this Hilbert space.
Thus the essential growth bound is at most zero. The standard
semigroup spectral decomposition outside the essential growth bound
\cite[Corollary~IV.2.11 and Theorem~V.3.1]{engel-nagel2000} makes all growing spectrum discrete of finite
algebraic multiplicity. Each such finite-dimensional rootspace is
rotation invariant, so the circle representation splits into integer
angular characters. The preceding transfer and exclusions leave
precisely $\lambda,\bar\lambda$, each simple. The same theorem gives
the first estimate in \eqref{eq:short-seed-complement}.
This statement includes the full $H^j_M$: splitting off a fixed radial
function of integral one merely adds a stationary direction.
Resolvents, hence Riesz projections, agree on intersections of
weighted spaces by uniqueness in a common weaker realization.

On $X_G$, the Green estimate
$|Rf(r)|\le C\|f\|_{X_G}\min(r,r^{-2})$ shows that $L_0$ is bounded.
After the unitary map $f\mapsto f/\sqrt G$, its interaction kernel is
\[
 -\frac i2\frac{k(r)}{\sqrt{G(r)}}
       \left(\frac{\min(r,s)}{\max(r,s)}\right)^2\sqrt{G(s)},
\]
which is Hilbert--Schmidt. Hence the same compact-perturbation and
spectral-decomposition argument gives the second estimate in
\eqref{eq:short-seed-complement}.

\emph{Actual adjoint, strain coupling, and diffusion domains.}
Unweighted selfadjointness of $R$ and $k=\Omega(x-2)$ give
\[
 L_0^*g=2i\Omega[g-R((x-2)g)].
\]
The explicit adjoint eigenvector is
\[
 h_X(r,\vartheta)=G(r)\frac{\overline{\phi(r)}}{
                    \Omega(r)-\bar c}e^{2i\vartheta},
 \qquad L_0^*h_X=\bar\lambda h_X.
\]
Indeed the conjugate Rayleigh equation gives
$R(k\bar\phi/(\Omega-\bar c))=\bar\phi$, with no homogeneous
ambiguity by its endpoints. Every derivative of $h_X$ is a polynomial
times $G$. The Fredholm operator $L_0-\lambda$ has index zero,
one-dimensional kernel, and adjoint kernel spanned by $h_X$.
Its pairing with $q$ cannot vanish, since otherwise $q$ would lie
in the range and produce a Jordan vector. This pairing is exactly
$J_q$ in \eqref{eq:short-seed-functional}. Setting
$q^\dagger=h_X/\overline{J_q}$
therefore gives that formula and
$\ell(H)=M_q/J_q\ne0$. The kernel $h/k=\phi/(\Omega-c)$ is smooth
at zero and $O(r^{-2})$ at infinity, so the same functional is bounded
on each polynomially weighted space. Its adjoint identity extends from
smooth data by density and gives the corresponding Riesz projection.

Finally, conjugation by $U_Gf=f/\sqrt G$ gives
\[
 U_GDU_G^{-1}=\Delta-|y|^2/16+3/2.
\]
The selfadjoint harmonic oscillator has Hermite levels
$1-N/2$; angular weight two consists exactly of $N=2+2m$,
$m\ge0$. This proves the domain and spectral assertion about $D$.
The Gaussian derivative bounds make $U_Gq$ and $U_Gh_X$ Schwartz,
so both belong to every oscillator power domain. With
$\ell(f)=\langle f,q^\dagger\rangle$, selfadjointness gives
$\ell(Df)=\langle f,Dq^\dagger\rangle$. Consequently
$PD=q\langle\,\cdot\,,Dq^\dagger\rangle$ and
$DP=(Dq)\ell$ are bounded.
\end{proof}

\begin{corollary}[Frozen heat profiles]\label{cor:short-seed-dilation}
Set
\[
 J_af(y)=a^{-2}f(y/\sqrt a),\qquad
 E_a=J_aE_0,\qquad V_a(y)=a^{-3/2}V_0(y/\sqrt a).
\]
Their Euler generators and unstable projections satisfy
\[
 \mathscr L_a=a^{-2}J_a\mathscr LJ_a^{-1},\qquad
 \mathscr P_a=J_a\mathscr PJ_a^{-1}.
\]
Thus the unstable eigenvalues are $a^{-2}\lambda,a^{-2}\bar\lambda$;
the complementary estimates are uniform for $1\le a\le A<\infty$,
and the projections depend continuously on $a$ in operator norm.
\end{corollary}
\begin{proof}
The generator identity is the Biot--Savart scaling identity.
$J_a,J_a^{-1}$ are uniformly bounded on $H^j_M$ on each compact
$a$ interval, so conjugating \eqref{eq:short-seed-complement} proves
the bounds. The finite-rank projection consists of the Gaussian
eigenvectors and the dual functional \eqref{eq:short-seed-functional};
their dilates are norm continuous by dominated convergence.
\end{proof}
